\subsection{Motivations: un problème hyper-sensible}
\label{part:hypersens}

L'article de Robert P. Taylor, R.~Luck et B.~K. Hedge~\cite{hypersensitive} met en lumière des problèmes hyper-sensibles, c'est-à-dire des problèmes pour lesquels une petite erreur dans les facteurs de vue peut engendrer une grande erreur lorsqu'elle est propagée au flux radiatif. En effet, leur étude de propagation d'incertitude montre une sensibilité extrême du flux de chaleur aux facteurs de vue. Dans le cas d'une géométrie pourtant simple d'un parallélépipède, traité dans~\cite{incro}, une précision de $5\,\%$ sur le flux de chaleur nécessite une précision de $10^{-4}$ sur les facteurs de vue.

Il devient donc nécessaire de maîtriser la précision des facteurs de vue calculés, et non uniquement de disposer d'une estimation ponctuelle : c'est cet enjeu qui motive le recours à des algorithmes de renforcement, présentés dans la section suivante.

\subsection{Algorithmes de renforcement}
\label{part:renforcement}
Les résultats obtenus par la méthode de lancer de $N$ rayons permettent une bonne approximation des facteurs de vue, avec une précision liée à l'estimateur en $\frac{1}{\sqrt{N}}$ caractéristique des méthodes de Monte Carlo. Des méthodes de correction~\cite{renforcement95}, qui reposent sur les propriétés géométriques des facteurs de vue, permettent d'améliorer la convergence de la matrice vers la matrice théorique. Même si nous disposons aujourd'hui d'une plus grande puissance de calcul qu'à l'époque où ces études ont été menées, l'enjeu reste toujours d'utiliser la méthode la moins coûteuse en temps de calcul, en particulier pour des géométries comportant un grand nombre de surfaces, où la matrice des facteurs de vue devient rapidement volumineuse.

L'algorithme de correction (\textit{enforcement algorithm}) permet de corriger la matrice de facteurs de vue obtenue par le calcul Monte Carlo afin qu'elle satisfasse les propriétés physiques de réciprocité et de fermeture présentées précédemment. Plusieurs stratégies de correction sont étudiées dans \cite{renforcement95}, notamment une approche directe et une méthode fondée sur l'optimisation au sens des moindres carrés. Les auteurs montrent que cette dernière offre les meilleurs résultats en termes de précision tout en limitant la variance introduite par la correction. Elle est donc retenue dans ce travail pour construire une matrice de facteurs de vue respectant au mieux les contraintes physiques du problème.

\subsubsection{Algorithme naïf}

L'algorithme naïf procède en deux étapes successives, appliquées à la matrice brute des facteurs de vue $F_{ij}$ issue du calcul Monte Carlo.

La première étape impose la relation de réciprocité $A_i F_{ij} = A_j F_{ji}$. Les termes hors diagonale étant a priori estimés indépendamment pour chaque paire $(i,j)$, rien ne garantit que $A_i F_{ij}$ et $A_j F_{ji}$ coïncident. On corrige alors chaque paire de termes symétriques par leur moyenne pondérée par les aires, de sorte que la matrice corrigée vérifie exactement la réciprocité :

$$\begin{bmatrix}
                ... & F_{12} & F_{13}\\
                ... & ... & F_{23} \\
                ... & ... & ... 
        \end{bmatrix}
        \xrightarrow{\text{réciprocité}}
        \begin{bmatrix}
                ... & F_{12} & F_{13}\\
                F_{21} & ... & F_{23} \\
                F_{31} & F_{32} & ... 
        \end{bmatrix}$$

La seconde étape impose la relation de fermeture $\sum_{j=1}^n F_{ij} = 1$ pour chaque surface $i$. Une fois la réciprocité assurée, chaque ligne de la matrice est normalisée de sorte que sa somme vaille exactement 1, ce qui permet également de faire apparaître les termes diagonaux $F_{ii}$ (non nuls dans le cas de surfaces non convexes) :

$$\begin{bmatrix}
                ... & F_{12} & F_{13}\\
                F_{21} & ... & F_{23} \\
                F_{31} & F_{32} & ... 
        \end{bmatrix}
        \xrightarrow{\text{fermeture}}
        \begin{bmatrix}
                F_{11} & F_{12} & F_{13}\\
                F_{21} & F_{22} & F_{23} \\
                F_{31} & F_{32} & F_{33} 
        \end{bmatrix}$$

Cette méthode a l'avantage de la simplicité et d'un faible coût de calcul. Elle ne garantit cependant pas la positivité de la matrice corrigée obtenue. En effet, pour $i \neq j$, la correction de réciprocité définit $\tilde{F}_{ij}$ et $\tilde{F}_{ji}$ par une moyenne pondérée par les aires :

\begin{equation}
    A_i \tilde{F}_{ij} = A_j \tilde{F}_{ji} = \frac{A_i F_{ij} + A_j F_{ji}}{A_i + A_j},
\end{equation}

puis la correction de fermeture impose, pour chaque surface $i$,

\begin{equation}
    \tilde{F}_{ii} = 1 - \sum_{j \neq i} \tilde{F}_{ij}.
\end{equation}

Rien ne garantit alors que $\tilde{F}_{ii}$ ainsi obtenu reste positif : si la somme $\sum_{j \neq i} \tilde{F}_{ij}$ dépasse l'unité, en raison des erreurs statistiques accumulées sur les termes hors diagonale lors de l'étape de réciprocité, le terme diagonal corrigé devient négatif, ce qui est physiquement impossible.


\subsubsection{Algorithme des moindres carrés}
\red{à rebosser d'ici le jour J}
La méthode des moindres carrés reformule la correction des facteurs de vue comme un problème d'optimisation sous contraintes. On cherche la matrice corrigée $F$ qui minimise l'écart quadratique pondéré à la matrice brute $F^0$ issue du calcul Monte Carlo :

\begin{equation}
    y = \sum_{i=1}^{n} \sum_{j=1}^{n} w_{ij} \left(F_{ij} - F_{ij}^0\right)^2
    \label{eq:moindre_carre}
\end{equation}

où les poids $w_{ij}$ permettent de pondérer chaque terme selon l'incertitude associée à son estimation Monte Carlo (les facteurs de vue estimés à partir d'un plus grand nombre de rayons pouvant par exemple se voir attribuer un poids plus élevé). Cette minimisation est soumise aux contraintes de fermeture et de réciprocité, qui s'écrivent sous forme d'égalités linéaires :

\begin{equation}
    \forall i, \quad \sum_{j=1}^{n} F_{ij} = 1 
    \label{eq:contrainte_fermeture}
\end{equation}

\begin{equation}
    \forall i \neq j, \quad A_j F_{ji} - A_i F_{ij} = 0
    \label{eq:contrainte_reciprocite}
\end{equation}

Une contrainte de positivité $F_{ij} \geq 0$ peut également être ajoutée si nécessaire, transformant le problème en un programme quadratique sous contrainte d'inégalité.

Plutôt que de résoudre ce problème par des méthodes générales de programmation non linéaire, \cite{renforcement95} propose une résolution reposant sur une interprétation géométrique du problème. Les facteurs de vue sont réorganisés sous la forme d'un vecteur colonne $f$ (obtenu en empilant les lignes de la matrice $F^0$), ce qui permet de réécrire les contraintes de fermeture et de réciprocité~\ref{eq:contraintes} sous la forme d'un système linéaire :

\begin{equation}
    R \cdot f = c
    \label{eq:systeme_contraintes}
\end{equation}

La solution de ce système peut être décomposée en deux composantes orthogonales : une solution particulière $f_{row}$, appartenant à l'espace ligne de $R$ et de norme minimale, donnée par

\begin{equation}
    f_{row} = R^T (R R^T)^{-1} c
    \label{eq:solution_particuliere}
\end{equation}

et une composante homogène appartenant au noyau de $R$, qui s'exprime comme une combinaison linéaire des vecteurs d'une base $N_b$ de ce noyau :

\begin{equation}
    f - f_{row} = N_b x
    \label{eq:solution_homogene}
\end{equation}

Les facteurs de vue bruts $f^0$ ne vérifiant pas exactement les contraintes, l'équation~\eqref{eq:solution_homogene} n'admet pas de solution exacte : le vecteur de poids $x$ est alors déterminé au sens des moindres carrés, en projetant l'écart $(f^0 - f_{row})$ sur le noyau de $R$ :

\begin{equation}
    x = (N_b^T N_b)^{-1} N_b^T (f^0 - f_{row})
    \label{eq:solution_x}
\end{equation}

Le vecteur des facteurs de vue corrigés $f$, reconstruit à partir des équations~\eqref{eq:solution_particuliere}, \eqref{eq:solution_homogene} et \eqref{eq:solution_x}, constitue alors la projection de $f^0$ sur l'espace des solutions vérifiant exactement les contraintes de fermeture et de réciprocité. Contrairement à l'algorithme naïf, cette approche traite l'ensemble des $n^2$ facteurs de vue de façon symétrique et cohérente, sans privilégier arbitrairement une partie de la matrice (comme le fait l'algorithme naïf en ne conservant que les termes au-dessus de la diagonale) ni recourir à un processus itératif de convergence incertaine (comme l'algorithme de Leersum). Cette rigueur a cependant un coût : le vecteur $f$ comportant $n^2$ composantes, la matrice $N_b^T N_b$ à inverser dans l'équation~\eqref{eq:solution_x} est de taille de l'ordre de $n^2 \times n^2$, ce qui conduit à une complexité de calcul évoluant en $\mathcal{O}(n^4)$. Cette méthode est ainsi vraisemblablement inadaptée à des cas 3D de grande taille, pour lesquels le nombre de surfaces $n$ peut être élevé.